%% 4-0 序章　ロボットの手は、どこにある？（第13回：組 V4-1。ページ制限なし。偶数ページ）
%% ① 長さと角度で場所が決まる（北東へ100 m。3-2 の座標とくらべる）② 順運動学と逆運動学（2-14 が「逆運動学とよぶ（4-0）」と指している）
%% ③ 3つの部の地図 ④ 通る路線と駅 ⑤ 新しい記号と、この巻の読み方（ルビがなくなる。「式」「答え」）
%% sin・cos は使うときに「これから学ぶしるし」（\bkivakore）をそえる。
\makeatletter
\gdef\bk@partlabel{第4巻}\gdef\bk@partname{形と向き}
%% ---- この章だけで使う部品（頭文字 \bkiva…） ----
% ロボットのうで  \bkivaarm[1ならつかむ手をかく]{(肩)}{上うでの長さ mm}{前うでの長さ mm}{肩の角}{ひじの角}{色}
%   (kivaS)・(kivaE)・(kivaH) に、肩・ひじ・手の位置の名前がつく
\newcommand\bkivaarm[7][0]{%
  \coordinate (kivaS) at #2;
  \coordinate (kivaE) at ($#2+(#5:#3)$);
  \coordinate (kivaH) at ($(kivaE)+({#5+(#6)}:#4)$);
  \draw[#7, line width=3.2pt, line cap=round] (kivaS) -- (kivaE) -- (kivaH);
  \ifnum#1=1
  \begin{scope}[shift={(kivaH)}, rotate={#5+(#6)}]
    \draw[bkInk, line width=1pt, line cap=round] (0.4,-1.7) -- (0.4,1.7) (0.4,-1.7) -- (2.4,-1.7) (0.4,1.7) -- (2.4,1.7);
  \end{scope}
  \fi
  \filldraw[fill=white, draw=#7!70!black, line width=0.8pt] (kivaS) circle (1.5);
  \filldraw[fill=white, draw=#7!70!black, line width=0.8pt] (kivaE) circle (1.2);}
% 方眼と軸  \bkivagrid{はしの x}{はしの y}{目もりの間 mm}{1目もりの値}{数字を書く間（目もり）}{横の軸の名前}{たての軸の名前}  ※原点は (0,0)
\newcommand\bkivagrid[7]{%
  \draw[bkLine, line width=0.3pt, step=#3] (0,0) grid (#1,#2);
  \draw[-{Stealth[length=3pt]}, bkInk, line width=0.6pt] (0,0) -- (#1+4,0) node[anchor=west, font=\sffamily\bfseries\footnotesize, inner sep=1pt] {#6};
  \draw[-{Stealth[length=3pt]}, bkInk, line width=0.6pt] (0,0) -- (0,#2+4) node[anchor=south, font=\sffamily\bfseries\footnotesize, inner sep=1pt] {#7};
  \pgfmathtruncatemacro\bkiva@n{floor(#1/(#3*#5)+0.01)}%
  \foreach \k in {1,...,\bkiva@n} {\pgfmathtruncatemacro\bkiva@v{\k*#4*#5}%
    \node[anchor=north, font=\sffamily\footnotesize, text=bkGray, inner sep=0pt] at ({\k*#3*#5},-1.6) {\bkiva@v};}
  \pgfmathtruncatemacro\bkiva@m{floor(#2/(#3*#5)+0.01)}%
  \foreach \k in {1,...,\bkiva@m} {\pgfmathtruncatemacro\bkiva@v{\k*#4*#5}%
    \node[anchor=east, font=\sffamily\footnotesize, text=bkGray, inner sep=0pt] at (-1.6,{\k*#3*#5}) {\bkiva@v};}
  \node[anchor=north east, font=\sffamily\footnotesize, text=bkGray, inner sep=0pt] at (-1.2,-1.6) {0};}
% 例題の札の形  \bkivapill{文字}
\newcommand\bkivapill[1]{\tikz[baseline=(s.base)]\node[fill=bkBlue, text=white, rounded corners=1.2mm,
    font=\sffamily\bfseries\small, inner xsep=4pt, inner ysep=2pt] (s) {#1};}
% これから学ぶしるし  \bkivakore
\newcommand\bkivakore{\tikz[baseline=(k.base)]\node[draw=bkGray, line width=0.6pt, fill=bkGray!10, rounded corners=1mm,
  inner xsep=1.2mm, inner ysep=0.6mm, font=\sffamily\bfseries\footnotesize, text=bkGray!60!black] (k) {これから学ぶ};}
% 式を、まん中に1行で置く  \bkivafline{…}
\newcommand\bkivafline[1]{\par\addvspace{3pt}{\centering #1\par}\addvspace{3pt}\noindent\ignorespaces}
% 箇条書きの1項目（・で始め、2行目は字下げ）  \bkivaitem{文}
\newcommand\bkivaitem[1]{\par\noindent\hangindent=1\zw\hangafter=1 ・#1\par}
% 路線の1行  \bkivaroute{y}{路線番号}{前の巻}{この巻の駅}{この先}
\newcommand\bkivaroute[5]{%
  \draw[bkRoute#2, line width=3pt, line cap=round] (10,#1) -- (132,#1);
  \node[inner sep=0pt] at (4,#1) {\bkrouteicon[7mm]{#2}};
  \fill[white] (24,#1) circle (2.2);\draw[bkRoute#2, line width=1.1pt] (24,#1) circle (2.2);
  \node[anchor=north, font=\sffamily\footnotesize, text=bkGray, inner sep=0pt] at (24,#1-3.8) {#3};
  \fill[white, rounded corners=1.5mm] (36,#1-4.4) rectangle (112,#1+4.4);
  \draw[bkRoute#2, line width=1.6pt, rounded corners=1.5mm] (36,#1-4.4) rectangle (112,#1+4.4);
  \node[font=\sffamily\footnotesize, text=bkInk] at (74,#1) {#4};
  \fill[white] (124,#1) circle (2.2);\draw[bkRoute#2, line width=1.1pt] (124,#1) circle (2.2);
  \node[anchor=north, font=\sffamily\footnotesize, text=bkGray, inner sep=0pt] at (124,#1-3.8) {#5};}
% 部の地図の1まい  \bkivapart{上のはしの y}{部の名前}{章}{絵}{中身}   はば132mm・高さ38mm
\newcommand\bkivapart[5]{%
  \draw[bkTheme!80!black, line width=1pt, fill=white, rounded corners=2.5mm] (0,#1) rectangle (132,#1-46);
  \node[anchor=north west, fill=bkTheme, text=bkOn, rounded corners=1.2mm, font=\sffamily\bfseries\small,
    inner xsep=2mm, inner ysep=1mm] at (3.5,#1-3.5) {#2};
  \node[anchor=north west, font=\sffamily\footnotesize, text=bkGray, inner sep=0pt] at (4.5,#1-11.5) {#3};
  \begin{scope}[shift={(22,#1-27)}]#4\end{scope}
  \node[anchor=north west, text width=78mm, align=left, font=\small, inner sep=0pt] at (51,#1-4.5) {#5};}
% 会話の1行  \bkivatalk{y}{sora|mina|delta}{名前}{せりふ}
\newcommand\bkivatalk[4]{%
  \bkChar{#2}{(7,#1-6.5)}{12mm}
  \node[rectangle callout, callout absolute pointer={(15.5,#1)}, anchor=west, draw=bkTheme!80!black, line width=0.8pt,
    fill=white, rounded corners=2.5mm, inner xsep=3mm, inner ysep=1.8mm, text width=103mm, align=left,
    font=\small] at (21,#1) {#4};
  \node[anchor=north, font=\sffamily\footnotesize, text=bkGray, inner sep=0pt] at (7,#1-7.5) {#3};}
% 読み方のカード  \bkivacard{(x,y)}{アイコン}{色}{見出し}{中身}  はば64mm・高さ64mm
\newcommand\bkivacard[5]{%
  \node[anchor=north west, draw=#3, line width=1pt, fill=white, rounded corners=2.5mm,
    minimum width=64mm, minimum height=48mm, inner sep=0pt] at #1 {};
  \node[anchor=north west, inner sep=0pt] (ci) at ($#1+(3.5,-3)$) {\bkicon[1.5em]{#2}{#3}};
  \node[anchor=west, font=\sffamily\bfseries\small, text=#3!85!black, inner sep=0pt] at ([xshift=2mm]ci.east) {#4};
  \node[anchor=north west, text width=57mm, align=left, font=\small, inner sep=0pt] at ($#1+(3.5,-12)$) {#5};}
% 見出しのアイコン（路線：線と2つの駅。黄緑の地なので線は濃い紺）
\newcommand\bk@icon@bkivaroute{%
  \draw[bkOn, line width=0.13*\bk@iu, line cap=round] (-0.62,0) -- (0.62,0);
  \filldraw[fill=white, draw=bkOn, line width=0.09*\bk@iu] (-0.36,0) circle (0.2) (0.36,0) circle (0.2);}
\makeatother

\bkchapter{4}{0}{ロボットの手は、\\どこにある？}{\bklevel{2}}{ロボットの手は、どこにある？}

\noindent 自動車の工場では、たくさんのロボットのうでが、部品をつかんだり、ねじをしめたり、溶接（ようせつ：金属をとかしてつなぐこと）をしたりしています。ロボットのうでは、人のうでと同じように、いくつかの「関節（かんせつ）」で曲がります。

\begin{bkfigure}
\begin{tikzpicture}[x=1mm, y=1mm]
  \fill[bkGray!25] (54,-1.5) rectangle (132,0);
  \draw[bkGray, line width=0.8pt] (54,0) -- (132,0);
  % 台座は床の線の上にじかに置き（下の角は丸めない）、その上の台に肩の関節をのせる（第14回 B1）
  \fill[bkGray!60] (59,0) [rounded corners=0.8mm] -- (59,6.5) -- (69,6.5) [sharp corners] -- (69,0) -- cycle;
  \fill[bkGray!80] (60.8,6.5) -- (67.2,6.5) -- (65.6,9) -- (62.4,9) -- cycle;
  \bkivaarm[1]{(64,9)}{28}{22}{55}{-85}{bkBlue}
  \node[anchor=west, font=\sffamily\small, inner sep=1pt] at (71,4.5) {肩（関節）};
  \node[anchor=south east, font=\sffamily\small, text=bkBlue!80!black, inner sep=1pt] at (70.5,22) {上うで};
  \node[anchor=south, font=\sffamily\small, inner sep=1pt] at (80.06,34.2) {ひじ（関節）};
  \node[anchor=south west, font=\sffamily\small, text=bkBlue!80!black, inner sep=1pt] at (91,28) {前うで};
  \node[anchor=west, font=\sffamily\small, inner sep=1pt] at (103.5,20.9) {手};
  \node[anchor=north west, text width=50mm, align=left, font=\small, inner sep=0pt] at (0,36)
    {ロボットのうでの部分は、人のうでと同じ名前でよびます。関節にはモーターが入っていて、うでを回す角度を決めます。うでの長さは変わりません。};
\end{tikzpicture}
\bkfigcaption{ロボットのうでの部分の名前（肩とひじの2つの関節で曲がるうで）}
\end{bkfigure}

ロボットを動かすコンピュータが、じかに決めているのは、関節のモーターを「何度回すか」だけです。手の位置は、だれも決めていません。それなのに、ロボットの手は、ベルトコンベアの上の小さな部品のま上に、ぴたりと止まります。

\begin{hatena}
\par\addvspace{4pt}\noindent\begin{minipage}[c]{0.56\linewidth}
うでの長さと、関節の角度がわかれば、手の位置は決まるかな？\par\vspace{3pt}
反対に、「手をここに置きたい」という場所から、関節の角度は決まるかな？
\end{minipage}\hfill
\begin{minipage}[c]{0.4\linewidth}\centering
\begin{tikzpicture}[x=1mm, y=1mm]
  % 床・台・ベルトコンベア・部品
  \fill[bkGray!25] (-2,-1.5) rectangle (48,0);
  \draw[bkGray, line width=0.8pt] (-2,0) -- (48,0);
  \fill[bkGray!60] (3,0) [rounded corners=0.8mm] -- (3,9.5) -- (13,9.5) [sharp corners] -- (13,0) -- cycle;
  \fill[bkGray!80] (4.8,9.5) -- (11.2,9.5) -- (9.6,12) -- (6.4,12) -- cycle;
  \fill[bkGray!45, rounded corners=1.5mm] (22,2.5) rectangle (48,6.5);
  \foreach \x in {24.5,29.5,34.5,39.5,44.5} {\fill[bkGray!70] (\x,4.5) circle (1.3);}
  \filldraw[fill=bkOrange!60, draw=bkOrange!80!black, line width=0.5pt] (31,6.5) rectangle (37,10.5);
  % うで（肩 (8,12)、上うで 23.86 mm・54.07°、前うで 20 mm・向き −60°。手は (32,14)）
  \bkivaarm[1]{(8,12)}{23.86}{20}{54.07}{-114.07}{bkBlue}
  \node[font=\sffamily\bfseries\Large, text=bkOrange!80!black] at (43,24) {?};
\end{tikzpicture}
\end{minipage}
\end{hatena}


この巻は、長さと角度を使って、「形」と「向き」を数で表す巻です。序章では、ロボットのうでを例にして、この巻で学ぶことの地図を作ります。

\begin{omoidasou}
\bkrecall{1-20}{1回転は360°、直角は90°。角の大きさは分度器ではかる。}
\bkrecall{3-2}{点 $(x,\ y)$ は、原点から右へ $x$、上へ $y$ の所。平面の上の場所は、2つの数で表せる。}
\bkrecall{2-14}{式を変形して、ほしい量をとり出す。ロボットの手の位置の式を、角度について解くことを「逆運動学」とよぶ、と予告した。}
\end{omoidasou}

\begin{yattemiyou}[紙のうでロボットを動かそう]
\noindent\begin{minipage}[c]{0.6\linewidth}
\textbf{用意するもの}\par
厚紙、割りピン2本、\qty{1}{cm} の方眼紙、分度器、定規
\begin{enumerate}[label={\bkcircled[bkGreen]{\arabic*}}, leftmargin=2.4\zw, itemsep=1pt, topsep=3pt]
\item 厚紙で、穴と穴のあいだが\qty{6}{cm}の上うでと\qty{4}{cm}の前うでを作り、割りピンでつなごう（ひじ）。方眼紙を厚紙にはり、上うでのはしを原点に割りピンでとめよう（肩）。
\item 分度器で、肩の角を30°、ひじの角を60°にしよう。手の先の穴は、どの点にあるかな？
\item 反対に、手の先の穴を点 $(6,\ 4)$ に置こう。置き方は何通りあるかな？
\item 手の先の穴を点 $(10,\ 3)$ に置けるかな？
\end{enumerate}
\end{minipage}\hfill
\begin{minipage}[c]{0.36\linewidth}\centering
\begin{tikzpicture}[x=1mm, y=1mm]
  \draw[bkLine, line width=0.3pt, step=4] (0,0) grid (36,32);
  \draw[bkInk, line width=0.6pt] (0,0) -- (36,0) (0,0) -- (0,32);
  % 上うで（6 cm）と前うで（4 cm）。1 cm を 4 mm でかく
  \begin{scope}[rotate=30]
    \filldraw[fill=bkYellow!55, draw=bkYellow!50!black, line width=0.5pt, rounded corners=1.6mm] (-2,-1.8) rectangle (26,1.8);
    \fill[white] (0,0) circle (0.7) (24,0) circle (0.7);
  \end{scope}
  \begin{scope}[shift={(30:24)}, rotate=90]
    \filldraw[fill=bkOrange!45, draw=bkOrange!60!black, line width=0.5pt, rounded corners=1.6mm] (-2,-1.8) rectangle (18,1.8);
    \fill[white] (16,0) circle (0.7);
  \end{scope}
  \fill[bkGray!70!black] (0,0) circle (0.9) (30:24) circle (0.9);
  \node[anchor=north, font=\sffamily\footnotesize, inner sep=0pt] at (0,-4.2) {肩};
  \node[anchor=west, font=\sffamily\footnotesize, inner sep=0pt] at ($(30:24)+(2.4,-1.2)$) {ひじ};
  \node[anchor=west, font=\sffamily\footnotesize, inner sep=0pt] at ($(30:24)+(3.5,16)$) {手};
\end{tikzpicture}
\end{minipage}
\end{yattemiyou}

\noindent ②では、手の先は点 $(5.2,\ 7.0)$ の近くに来たはずです。③と④の答えは、このあと「ロボットのうでの2つの問題」で考えます。

%% ---------------- ① 長さと角度で場所が決まる ----------------
\bksection{doguu}{長さと角度で、場所が決まる}

\bksubsection{1}{「北東へ100 m」}

\label{bk:4-0:basho}遠足の集合場所を、友だちに伝えます。2つの伝え方をくらべてみましょう。

\begin{bkcenter}
\begin{bktable}{colspec={Q[l,m,wd=30mm] Q[l,m,wd=88mm]}, rowsep=2pt, row{1}={font=\sffamily\bfseries\small}}
伝え方 & 伝えることば \\
ア　たてと横 & 「校門から、東へ約\qty{71}{m}、北へ約\qty{71}{m}行った所」 \\
イ　長さと向き & 「校門から、北東へ\qty{100}{m}行った所」 \\
\end{bktable}
\bktabcaption{同じ場所を、2通りに伝える}
\end{bkcenter}

アは、3-2の座標と同じ伝え方です。東を $x$ 軸の向き、北を $y$ 軸の向きにすると、集合場所は点 $(71,\ 71)$ です。イは、「まっすぐ行く長さ」と「向き」の2つで伝えています。北東は、東と北のちょうどまん中の向きです。東から北へ、直角の半分の45°だけ回した向きです。

\par\addvspace{4pt}\noindent\begin{minipage}[c]{0.5\linewidth}\centering
\begin{tikzpicture}[x=1mm, y=1mm]
  % 1 m を 0.4 mm でかく。方眼は 10 m ごと
  \bkivagrid{44}{44}{4}{10}{2}{東}{北}
  \draw[bkRed, line width=0.7pt, densely dashed] (28.28,28.28) -- (28.28,0) (28.28,28.28) -- (0,28.28);
  \draw[-{Stealth[length=5pt]}, bkBlue, line width=1.3pt] (0,0) -- (28.28,28.28);
  \fill[bkOrange] (28.28,28.28) circle (1.3);
  \node[anchor=south west, font=\sffamily\footnotesize, inner sep=1pt, fill=white] at (29.3,29.3) {集合場所};
  \draw[bkInk, line width=0.6pt] (9,0) arc (0:45:9);
  \node[font=\sffamily\footnotesize, fill=white, inner sep=0.8pt] at (22.5:14) {45°};
  \node[anchor=south east, font=\sffamily\footnotesize, text=bkBlue!80!black, fill=white, inner sep=1pt, rotate=45] at (14,16.5) {\qty{100}{m}};
\end{tikzpicture}
\end{minipage}\hfill
\begin{minipage}[c]{0.47\linewidth}
方眼の1目もりを\qty{10}{m}にして（目もりの数字の単位は m）、校門を原点にかきました。\par\vspace{3pt}
青い矢印は、東から45°の向きへ\qty{100}{m}（10目もり）の長さです。定規で\qty{100}{m}分の長さをはかってかきます。\par\vspace{3pt}
矢印の先から、赤い点線を軸までおろして、目もりを読みます。東へも北へも、7目もりと少し（約\qty{71}{m}）です。
\end{minipage}

\par\vspace{4pt}
長さと向きがわかれば、図をかいて、たてと横の長さを読みとれます。反対に、たてと横の長さがわかれば、図をかいて、矢印の長さと向きをはかれます。アとイは、同じ場所の、2通りの言い方です。

\bksubsection{2}{向きを変えると、たてと横はどう変わる？}

向きの角を変えて、同じように図をかいて読みとってみます。角は、東から左回り（北のほう）にはかります。

\begin{bkfigure}
\begin{tikzpicture}[x=1mm, y=1mm]
  \bkivagrid{44}{44}{4}{10}{2}{東}{北}
  \foreach \a/\c in {30/bkOrange, 45/bkBlue, 60/bkPurple} {
    \draw[-{Stealth[length=4pt]}, \c, line width=1.1pt] (0,0) -- (\a:40);}
  \draw[bkOrange!80!black, line width=0.6pt, densely dashed] (30:40) -- (34.64,0);
  \node[anchor=west, font=\sffamily\footnotesize, text=bkOrange!80!black, inner sep=1pt] at (30:41.5) {30°};
  \node[anchor=south west, font=\sffamily\footnotesize, text=bkBlue!80!black, inner sep=1pt] at (45:40.5) {45°};
  \node[anchor=south west, font=\sffamily\footnotesize, text=bkPurple!80!black, inner sep=1pt] at (60:40.5) {60°};
  % 右の説明
  \node[anchor=north west, text width=66mm, align=left, font=\small] at (60,44) {%
  長さはどれも\qty{100}{m}（10目もり）。\\[3pt]
  オレンジの30°の矢印は、横に約8.7目もり、たてに5目もりの点に着く。\\[3pt]
  だから「東へ約\qty{87}{m}、北へ\qty{50}{m}」。\\[3pt]
  角を大きくするほど、矢印は北へ回り、北への長さがふえて、東への長さがへる。};
\end{tikzpicture}
\bkfigcaption{東から30°・45°・60°の向きへ、\qty{100}{m}ずつ進んだ所。0°はま東、90°はま北。}
\end{bkfigure}

\begin{bkcenter}
\begin{bktable}{colspec={Q[l,m,wd=40mm] *{5}{Q[c,m,wd=14mm]}}, rowsep=1.5pt}
向き（東から左回りの角） & 0° & 30° & 45° & 60° & 90° \\
東へ（m） & 100 & 約87 & 約71 & 50 & 0 \\
北へ（m） & 0 & 50 & 約71 & 約87 & 100 \\
\end{bktable}
\bktabcaption{\qty{100}{m}進んだときの、東への長さと北への長さ（図から読みとった値）}
\end{bkcenter}

表から、次のことがわかります。
\begin{itemize}[leftmargin=1.2\zw, itemsep=2pt, topsep=2pt, label={\color{bkTheme!80!black}●}]
\item 角が決まると、東への長さと北への長さが1つに決まる。
\item 30°と60°では、東と北の数が入れかわる。30°の矢印と60°の矢印は、45°の線をはさんで、鏡に映したような形（線対称、4-1）だから。
\item もし\qty{200}{m}進むなら、同じ形のまま2倍に大きくなるので、東へも北へも2倍になる（にた形、4-3）。
\end{itemize}

角と長さから、たてと横の長さを、図をかかずに計算で求める道具が、$\cos$（コサイン）と $\sin$（サイン）です。\bkivakore\enspace
この印のついたことばや式は、これから学ぶものです。いまはわからなくてもだいじょうぶです。4-6で、直角三角形の辺の比として作り、「長さと向き」を「たてと横」に直す道具として使います。

\begin{taisetsu}
平面の上の場所は、2つの数で決まる。\par
「たての長さと横の長さ（座標）」でも、「まっすぐの長さと向きの角」でもよい。
\end{taisetsu}

「長さと角」で場所を表す方法は、レーダーや、ロボットのうでのように「回るもの」で使うと便利です。この方法は、4-2 で\textbf{極座標}（きょくざひょう）として学びます。極座標を、$\cos$・$\sin$ を使ってたてと横（座標）に直す方法は 4-6 と 4-8 で、反対に、たてと横から長さと角を求める方法や、空間への広げ方は 4-15 で学びます。

%% ---------------- ② 順運動学と逆運動学 ----------------
\bksection{doguu}{ロボットのうでの2つの問題}

\bksubsection{3}{角度から手の位置を求める}

\label{bk:4-0:jun}いよいよ、ロボットのうでです。肩からひじまでの上うでが\qty{30}{cm}、ひじから手までの前うでが\qty{20}{cm}のロボットを考えます（やってみようの紙のうでは、この5分の1の大きさです）。肩を原点に置き、長さは cm ではかります。

角は、次のように決めます。
\begin{itemize}[leftmargin=1.2\zw, itemsep=2pt, topsep=2pt, label={\color{bkTheme!80!black}●}]
\item \textbf{肩の角}：$x$ 軸の向きから、上うでがどれだけ左回りに回ったか。
\item \textbf{ひじの角}：上うでの向き（点線）から、前うでがどれだけ左回りに曲がったか。
\end{itemize}

\begin{bkfigure}
\begin{tikzpicture}[x=1mm, y=1mm]
  % 1 cm を 0.7 mm でかく。方眼は 10 cm ごと
  \foreach \sx/\ta/\tb/\c/\lab in {0/30/60/bkBlue/A, 46/60/60/bkOrange/B, 92/30/0/bkPurple/C} {
    \begin{scope}[shift={(\sx,0)}]
      {\draw[bkLine, line width=0.3pt, step=7] (0,0) grid (35,35);
      \draw[bkInk, line width=0.6pt] (0,0) -- (37,0) (0,0) -- (0,37);
      \foreach \k in {1,...,5} {\pgfmathtruncatemacro\v{\k*10}\node[anchor=north, font=\sffamily\footnotesize, text=bkGray, inner sep=0pt] at ({\k*7},-1.6) {\v};
        \node[anchor=east, font=\sffamily\footnotesize, text=bkGray, inner sep=0pt] at (-1.4,{\k*7}) {\v};}
      \bkivaarm{(0,0)}{21}{14}{\ta}{\tb}{\c}
      \draw[bkGray, line width=0.5pt, densely dashed] (kivaE) -- ($(kivaE)+(\ta:7)$);
      \draw[bkInk, line width=0.5pt] (5,0) arc (0:\ta:5);
      \node[font=\sffamily\footnotesize, inner sep=0.5pt] at ({\ta/2}:11) {\ta°};
      \ifnum\tb>0
        \draw[bkInk, line width=0.5pt] ($(kivaE)+(\ta:4)$) arc (\ta:{\ta+\tb}:4);
        \node[font=\sffamily\footnotesize, inner sep=0.5pt] at ($(kivaE)+({\ta+\tb/2}:10)$) {\tb°};
      \else
        \node[anchor=north west, font=\sffamily\footnotesize, inner sep=0.5pt] at ($(kivaE)+(-60:3)$) {0°};
      \fi
      \fill[bkInk] (kivaH) circle (1.1);
      \node[anchor=south, font=\sffamily\bfseries\small, text=\c!80!black] at (17.5,38) {\lab};}
    \end{scope}}
\end{tikzpicture}
\bkfigcaption{A・B・Cの3つの形。黒い点が手の位置。方眼の1目もりは\qty{10}{cm}。}
\label{bk:4-0:fig-abc}
\end{bkfigure}

\begin{bkcenter}
\begin{bktable}{colspec={Q[c,m,wd=12mm] Q[c,m,wd=20mm] Q[c,m,wd=22mm] Q[c,m,wd=34mm] Q[c,m,wd=26mm]}, rowsep=1.5pt, row{1}={font=\sffamily\small}}
形 & 肩の角 & ひじの角 & 前うでの向きの角 & 手の位置 \\
A & 30° & 60° & $30+60=90$° & 約 $(26,\ 35)$ \\
B & 60° & 60° & $60+60=120$° & 約 $(5,\ 43)$ \\
C & 30° & 0° & $30+0=30$° & 約 $(43,\ 25)$ \\
\end{bktable}
\bktabcaption{角度を決めると、手の位置が1つに決まる（図の目もりを読んだ値）}
\end{bkcenter}

手の位置は、①の「長さと向きで場所が決まる」を2回くり返すと決まります。
\begin{enumerate}[label={\bkcircled[bkDeep]{\arabic*}}, leftmargin=2.4\zw, itemsep=2pt, topsep=2pt]
\item 肩の角で上うでの向きが決まる。肩から\qty{30}{cm}その向きへ進んだ所が、ひじ。
\item 前うでの向きは、肩の角とひじの角をたした角になる（上うでの向きから、さらにひじの角だけ回すから）。
\item ひじから\qty{20}{cm}その向きへ進んだ所が、手。
\end{enumerate}

たとえばAでは、前うでの向きは $30+60=90$ で90°、つまりま上です。だから、手はひじのま上にあります。図でも、Aの前うではまっすぐ上を向いています。Bでは $60+60=120$ で120°なので、前うでは左ななめ上を向き、手はひじより左に来ます。Cは、ひじを曲げていない（0°）ので、うで全体が30°の向きにまっすぐのびています。①の30°の矢印と同じ形で、長さは $30+20=50$ で\qty{50}{cm}です。

角度を決めれば、手の位置はいつも1つに決まります。この計算を、ロボットの分野では\textbf{順運動学}（じゅんうんどうがく）といいます。「角度 → 位置」の順の向きに計算する、という意味です。\bkidx{じゅんうんどうがく}{順運動学}

\begin{seikaku}
\textbf{計算で書くと}\par
肩の角を $\theta_1$、ひじの角を $\theta_2$ と書くと（$\theta$ はシータと読むギリシャ文字で、角の大きさによく使う。4-2）、手の位置 $(x,\ y)$ は次の式で計算できます。\enspace\bkivakore
\[ x=30\cos\theta_1+20\cos(\theta_1+\theta_2),\qquad y=30\sin\theta_1+20\sin(\theta_1+\theta_2) \]
1つめの項が「肩からひじまで」、2つめの項が「ひじから手まで」です。$\theta_1+\theta_2$ が、②の「前うでの向きの角」です。この式は4-7で作ります。
\end{seikaku}

\bksubsection{4}{手を置きたい場所から、角度を求める}

\label{bk:4-0:gyaku}こんどは反対です。Aの手の位置（約 $(26,\ 35)$）に、ほかの角度で手を置けないでしょうか。

\par\addvspace{4pt}\noindent\begin{minipage}[c]{0.5\linewidth}\centering
\begin{tikzpicture}[x=1mm, y=1mm]
  % 1 cm を 1.1 mm でかく
  \draw[bkLine, line width=0.3pt, step=11] (0,0) grid (55,55);
  \draw[bkInk, line width=0.6pt] (0,0) -- (57,0) (0,0) -- (0,57);
  \foreach \k in {1,...,5} {\pgfmathtruncatemacro\v{\k*10}\node[anchor=north, font=\sffamily\footnotesize, text=bkGray, inner sep=0pt] at ({\k*11},-1.6) {\v};
    \node[anchor=east, font=\sffamily\footnotesize, text=bkGray, inner sep=0pt] at (-1.6,{\k*11}) {\v};}
  \draw[bkGray, line width=0.6pt, dash pattern=on 2pt off 1.5pt] (0,0) -- (28.58,38.5);
  \bkivaarm{(0,0)}{33}{22}{30}{60}{bkBlue}
  \bkivaarm{(0,0)}{33}{22}{76.79}{-60}{bkOrange}
  \fill[bkInk] (28.58,38.5) circle (1.3);
  \node[anchor=west, font=\sffamily\footnotesize, fill=white, inner sep=1pt] at (31,39.5) {手を置く点};
  \node[anchor=north west, font=\sffamily\footnotesize, text=bkBlue!80!black, fill=white, inner sep=1pt] at (29.5,14.5) {ひじが下（A）};
  \node[anchor=south, font=\sffamily\footnotesize, text=bkOrange!80!black, fill=white, inner sep=1pt] at (16,39.5) {ひじが上};
\end{tikzpicture}
\end{minipage}\hfill
\begin{minipage}[c]{0.47\linewidth}
同じ点に、2通りの形で手が届きます。\par\vspace{3pt}
\textcolor{bkBlue!80!black}{青}はAの形（肩30°、ひじ60°）です。\par\vspace{3pt}
\textcolor{bkOrange!80!black}{オレンジ}は、肩を約77°まで上げて、ひじを反対向き（右回り）に60°曲げた形です。右回りは、角にマイナスをつけて $-60$° と書きます（2-2：マイナスは反対の向き）。\par\vspace{3pt}
2つの形は、肩と手を結ぶ点線をはさんで、鏡に映したような形（線対称、4-1）です。
\end{minipage}

\par\vspace{4pt}
やってみようの③も同じです。紙のうでの手を $(6,\ 4)$ に置くには、「肩0°・ひじ90°」と「肩約67°・ひじ $-90$°」の2通りがあります。

このように、「手を置きたい位置 → 角度」の向きに計算することを、\textbf{逆運動学}（ぎゃくうんどうがく）といいます。2-14の「ロボットのうで」で予告した名前です。順運動学の式を、角度について解く計算です。\bkidx{ぎゃくうんどうがく}{逆運動学}

\begin{bkfigure}
\begin{tikzpicture}[x=1mm, y=1mm,
  q/.style={draw=bkTheme!80!black, line width=1.2pt, fill=white, rounded corners=2.5mm, minimum width=36mm, minimum height=16mm, align=center, font=\sffamily\small}]
  \node[q] (k) at (0,0) {\textbf{関節の角度}\\[1pt]{\footnotesize 肩の角・ひじの角}};
  \node[q] (p) at (84,0) {\textbf{手の位置}\\[1pt]{\footnotesize $(x,\ y)$}};
  \draw[bk arrow, bkBlue] ([yshift=3.5mm]k.east) -- node[above, font=\sffamily\bfseries\small, text=bkBlue!80!black] {順運動学} ([yshift=3.5mm]p.west);
  \draw[bk arrow, bkOrange] ([yshift=-3.5mm]p.west) -- node[below, font=\sffamily\bfseries\small, text=bkOrange!80!black] {逆運動学} ([yshift=-3.5mm]k.east);
\end{tikzpicture}
\bkfigcaption{順運動学は右へ、逆運動学は左へ。2つは逆向きの計算。}
\end{bkfigure}


では、手をどこにでも置けるのでしょうか。やってみようの④の点 $(10,\ 3)$ は、紙のうで（\qty{6}{cm}と\qty{4}{cm}）では届きません。うでをまっすぐにのばしても、届くのは肩から $6+4=10$ で\qty{10}{cm}までです。点 $(10,\ 3)$ は、肩から\qty{10}{cm}より少し遠いのです（どれだけ遠いかは、4-4の三平方の定理で計算できます）。

\par\addvspace{4pt}\noindent\begin{minipage}[c]{0.38\linewidth}\centering
\begin{tikzpicture}[x=0.82mm, y=0.82mm]
  % 本物のうで（30 cm と 20 cm）。1 cm を 0.55 mm でかく
  \fill[bkGreen!22, even odd rule] (0,0) circle (27.5) (0,0) circle (5.5);
  \draw[bkGray, line width=0.4pt] (-29.5,0) -- (29.5,0) (0,-29.5) -- (0,29.5);
  \draw[bkGreen!60!black, line width=0.9pt] (0,0) circle (27.5);
  \draw[bkGreen!60!black, line width=0.9pt] (0,0) circle (5.5);
  \fill[bkInk] (0,0) circle (1);
  \node[font=\sffamily\bfseries\footnotesize, fill=bkGreen!22, inner sep=0.8pt] at (13,-14) {届く};
  \draw[bkRed, line width=1pt] (19.5,20.5) -- (22.5,23.5) (19.5,23.5) -- (22.5,20.5);
\end{tikzpicture}
\end{minipage}\hfill
\begin{minipage}[c]{0.59\linewidth}
本物のうで（\qty{30}{cm}と\qty{20}{cm}）で、手が届く所を\textcolor{bkGreen!50!black}{緑}でぬりました。まん中の黒い点が肩です。外側の円は半径\qty{50}{cm}、内側の円は半径\qty{10}{cm}です。\par\vspace{3pt}
・肩から\qty{50}{cm}（$30+20$）より遠い所（赤い×など）には届かない。\par
・肩から\qty{10}{cm}（$30-20$）より近い所にも届かない。前うでを上うでにいちばん近くまで折りたたんでも、手は肩から\qty{10}{cm}はなれているから。\par
・緑の所なら、答えは2通り。外側の円の上（うでをまっすぐのばす）と、内側の円の上（いちばん折りたたむ）では、答えは1通り。
\end{minipage}

\par\vspace{4pt}
\begin{bkcenter}
\begin{bktable}{colspec={Q[l,m,wd=33mm] Q[l,m,wd=33mm] Q[l,m,wd=54mm]}, colsep=4pt, rowsep=2pt, row{1}={font=\sffamily\bfseries\small}}
 & 順運動学 & 逆運動学 \\
わかっていること & 関節の角度 & 手を置きたい位置 \\
求めること & 手の位置 & 関節の角度 \\
答えの数 & いつも1つ & 2つ・1つ・0のことがある \\
計算のしかた & 式に入れるだけ & 角度について解く（4-7） \\
\end{bktable}
\bktabcaption{順運動学と逆運動学のくらべ}
\end{bkcenter}

\bkhatenaref の答え：うでの長さと関節の角度がわかれば、手の位置は1つに決まります（順運動学）。反対に、手を置きたい場所から角度を求めると（逆運動学）、答えが2つあったり、届かなくて答えがなかったりします。逆の向きの計算のほうが、ずっとむずかしいのです。

\bkfukidashi{mina}{わたしもコップに手をのばすとき、ひじを下にしているな。ひじを上にしても、同じ所に届くんだね。}

\begin{taisetsu}
\textbf{順運動学}：角度 → 手の位置。答えはいつも1つ。\par
\textbf{逆運動学}：手の位置 → 角度。答えは2つ・1つ・0のことがある。
\end{taisetsu}

\begin{butsuri}
\par\addvspace{3pt}\textbf{ショベルカーの画面}\par\noindent
ショベルカーも、うでの3つの関節の角度をセンサーではかり、バケットの先の深さを順運動学で計算して、運転席の画面に表示します。

\par\addvspace{3pt}\textbf{人のうでも、逆運動学を解いている}\par\noindent
コップに手をのばすとき、わたしたちの脳は、目で見たコップの位置から、肩とひじを曲げる角度を決めています。
\end{butsuri}

\begin{butsuri}
\par\addvspace{3pt}\textbf{工場のロボットは、関節の角度をはかっている}\par\vspace{3pt}\noindent
\begin{minipage}[c]{0.66\linewidth}
工場のロボットのうでには、ふつう6つの関節があります。どの関節のモーターにも、回った角度をはかるセンサー（エンコーダー）がついています。細いすき間がたくさん並んだ円板がモーターといっしょに回り、すき間の列をはさんだ光のセンサー（半導体でできている）が、光が通った回数を数えます。その数から、回った角度がわかります。制御装置は、はかった角度から順運動学で手の位置を、1秒に1000回ほども計算しています。
\end{minipage}\hfill
\begin{minipage}[c]{0.3\linewidth}\centering
\begin{tikzpicture}[x=1mm, y=1mm]
  % エンコーダーの円板。右のすき間の列を、光の出口とセンサー（紺）ではさむ
  \fill[bkGray!30] (0,0) circle (13);
  \foreach \a in {0,15,...,345} {\fill[white] (\a:9.5) -- (\a+5:9.5) -- (\a+5:12) -- (\a:12) -- cycle;}
  \draw[bkGray!80!black, line width=0.6pt] (0,0) circle (13);
  \fill[bkGray!70!black] (0,0) circle (2);
  \fill[bkNavy, opacity=0.85, rounded corners=0.6mm] (8.6,-2.2) rectangle (15.5,2.2);
  \draw[bkNavy, line width=0.6pt] (15.5,0) -- (17,0) -- (17,-14);
  \node[anchor=north, font=\sffamily\footnotesize, align=center, inner sep=1pt] at (17,-14.5) {光と\\センサー};
  \draw[-{Stealth[length=3pt]}, bkOrange, line width=0.8pt] (-6,15) arc (115:65:14);
  \node[anchor=south, font=\sffamily\footnotesize, inner sep=1pt] at (0,17) {モーターと回る};
\end{tikzpicture}
\end{minipage}
\end{butsuri}

ロボットのうでについての、ほかの問題も、この巻のいろいろな章で、少しずつ解けるようになります。

\begin{bkcenter}
\begin{bktable}{colspec={Q[l,m,wd=78mm] Q[l,m,wd=42mm]}, rowsep=1.5pt, row{1}={font=\sffamily\bfseries\small}}
ロボットのうでについて知りたいこと & 解ける章 \\
肩から手までの、まっすぐな距離 & 4-4（三平方の定理） \\
関節が1秒に回る角度と、手の速さ & 4-2・4-8 \\
角度から手の位置を計算する式（順運動学） & 4-6・4-7 \\
手の位置から、ひじの角を計算する（逆運動学） & 4-7（余弦定理） \\
手が動く速さの向き、手が物をおす力の向き & 4-10・4-11 \\
モーターが関節を回す力（トルク） & 4-14（外積） \\
\end{bktable}
\bktabcaption{ロボットのうでの問題と、それが解ける章}
\end{bkcenter}

%% ---------------- ③ 3つの部の地図 ----------------
\bksection{goal}{この巻でできるようになること}

\label{bk:4-0:chizu}第4巻は、3つの部に分かれています。どの部も、①と②で見た「長さと角度で、形と向きを数で表す」ことを、少しずつくわしくしていきます。

\begin{bkfigure}
\begin{tikzpicture}[x=1mm, y=1mm]
  \bkivapart{0}{第1部\enspace 図形}{4-1〜4-5}%
    {\draw[bkBlue, line width=1.1pt, line join=round] (-12,-7) -- (12,-7) -- (3,9) -- cycle;
     \draw[bkInk, line width=0.5pt] (-7,-7) arc (0:33.7:5) (7,-7) arc (180:119.4:5);}
    {\bkivaitem{きまり（定義・公理）から、証明で考えを積み上げる。}\bkivaitem{角を「長さ」（ラジアン）ではかる。場所を「距離と向き」（極座標）で言う。にた形の比と三平方の定理で、長さを計算する。}\bkivaitem{面積や体積を、細かく分けてたして求める。}}
  \bkivapart{-50}{第2部\enspace 三角比と三角関数}{4-6〜4-9}%
    {\draw[bkInk, line width=0.6pt] (-10,0) circle (7);
     \draw[bkInk, line width=0.5pt] (-10,0) -- ({-10+7*cos(50)},{7*sin(50)});
     \fill[bkOrange] ({-10+7*cos(50)},{7*sin(50)}) circle (1.1);
     \draw[bkBlue, line width=1.1pt, domain=0:360, samples=40] plot ({0+\x/360*16}, {7*sin(\x+50)});}
    {\bkivaitem{距離と向きを、たてと横に直す（$\sin$・$\cos$）。坂の角度から、重さを2つの向きの力に分ける。}\bkivaitem{行けない場所までの距離や、ロボットのひじの角を計算する。}\bkivaitem{回る点の影から波を作り、2つの波をたし合わせる。}}
  \bkivapart{-100}{第3部\enspace ベクトル}{4-10〜4-15}%
    {\draw[-{Stealth[length=4pt]}, bkBlue, line width=1.2pt] (-12,-7) -- (6,-3);
     \draw[-{Stealth[length=4pt]}, bkOrange, line width=1.2pt] (-12,-7) -- (-6,8);
     \draw[bkGray, line width=0.5pt, dash pattern=on 1.5pt off 1.5pt] (6,-3) -- (12,12) -- (-6,8);
     \draw[-{Stealth[length=4.5pt]}, bkPurple, line width=1.5pt] (-12,-7) -- (12,12);}
    {\bkivaitem{力や速さを、向きと大きさをもつ矢印（ベクトル）で表し、たしたり分けたりする。}\bkivaitem{仕事（内積）と、回す力（外積）を計算する。}\bkivaitem{空間の点を、3つの数や、距離と角で表す。}}
\end{tikzpicture}
\bkfigcaption{第4巻の3つの部と、できるようになること}
\end{bkfigure}



%% ---------------- ④ この巻が通る路線と駅 ----------------
\bksection{bkivaroute}{この巻が通る路線と駅}

\label{bk:4-0:rosen}この本の章は、7本の路線の駅です。同じ路線の駅は、1つの考え方でつながっています。第4巻を通るのは、次の3本です。

\begin{bkfigure}
\begin{tikzpicture}[x=1mm, y=1mm]
  \bkivaroute{28}{2}{2-3}{{\bfseries 4-5}}{5-10}
  \bkivaroute{14}{4}{1-20}{{\bfseries 4-2}・4-6・4-7・{\bfseries 4-8}・4-9}{5-15}
  \bkivaroute{0}{6}{3-2}{4-1・4-3・4-4・{\bfseries 4-10}〜4-15}{7-7}
  \node[font=\sffamily\bfseries\small, text=bkGray] at (24,38) {前の巻};
  \node[font=\sffamily\bfseries\small, text=bkDeep] at (74,38) {この巻（第4巻）};
  \node[font=\sffamily\bfseries\small, text=bkGray] at (124,38) {この先};
\end{tikzpicture}
\bkfigcaption{第4巻を通る3本の路線。四角の中がこの巻の駅（太字は大きな駅）。}
\end{bkfigure}

\begin{bkcenter}
\begin{bktable}{colspec={Q[c,m,wd=7mm] Q[l,m,wd=20mm] Q[l,m,wd=92mm]}, rowsep=0.8pt, cells={font=\small}, row{1}={font=\sffamily\bfseries\small}}
 & 路線 & 大きな駅のつながり \\
\bkrouteicon{2} & たし合わせ & かけ算の面積（2-3）→ 細かく分けてたす面積と体積（4-5）→ 積分（5-10） \\
\bkrouteicon{4} & 回転と波 & 角の大きさ（1-20）→ ラジアン（4-2）・三角関数（4-8）\mbox{→ $\sin$・$\cos$ の微分（5-15）} \\
\bkrouteicon{6} & 矢印 & 座標（3-2）→ ベクトル（4-10）→ 3次元の回転（7-7）・ベクトル場（8-5） \\
\end{bktable}
\bktabcaption{3本の路線の、この巻の大きな駅}
\end{bkcenter}

序章のロボットのうでは、回転と波の路線（\bkrouteicon[0.9em]{4}）の駅です。関節がくるくる回ると、手は円をえがき、その影は波のように行ったり来たりします（4-8）。矢印の路線（\bkrouteicon[0.9em]{6}）は、この巻でいちばんたくさんの駅に止まります。点・線・角（4-1）から始まり、向きと大きさをもつ矢印（ベクトル、4-10）へと進みます。

\begin{column}{なぜ「形と向き」なの？}
第3巻では、時間がたつと量がどう変わるかを、グラフと式で表しました。第4巻では、「どこにあるか」「どっちを向いているか」を数で表します。ものの形（図形）、回る角度（三角関数）、力や速さの向き（ベクトル）です。物理のきまりの多くは、「どの向きに、どれだけの力がはたらくか」で書かれています。だから、第4巻の道具は、このあと学ぶ物理（力学）の土台になります。
\par\vspace{3pt}
{\centering
\begin{tikzpicture}[x=0.88mm, y=0.88mm]
  % 図形：橋のトラス
  \begin{scope}[shift={(0,0)}]
    \fill[bkGray!25] (-2,-2) rectangle (36,0);
    \draw[bkBlue, line width=1pt, line join=round] (0,0) -- (34,0) (5,10) -- (29,10)
      (0,0) -- (5,10) -- (11,0) -- (17,10) -- (23,0) -- (29,10) -- (34,0);
    \node[anchor=north, font=\sffamily\footnotesize, inner sep=0pt] at (17,-4) {図形：橋の三角形（4-1）};
  \end{scope}
  % 三角関数：観覧車
  \begin{scope}[shift={(58,7)}]
    \draw[bkInk!60, line width=0.8pt] (0,0) circle (9);
    \draw[bkInk!60, line width=0.7pt] (0,0) -- (-5,-11) (0,0) -- (5,-11);
    \foreach \a in {0,60,...,300} {\draw[bkInk!35, line width=0.5pt] (0,0) -- (\a:9); \fill[bkOrange] (\a:9) circle (1.2);}
    \draw[-{Stealth[length=3pt]}, bkRed, line width=0.8pt] (30:12) arc (30:80:12);
    \node[anchor=north, font=\sffamily\footnotesize, inner sep=0pt] at (0,-15) {回転：観覧車（4-8）};
  \end{scope}
  % ベクトル：2人で引く荷物
  \begin{scope}[shift={(96,2)}]
    \filldraw[fill=bkOrange!45, draw=bkOrange!70!black, line width=0.6pt] (0,0) rectangle (8,6);
    \draw[-{Stealth[length=4pt]}, bkBlue, line width=1.1pt] (8,4) -- (24,10);
    \draw[-{Stealth[length=4pt]}, bkBlue, line width=1.1pt] (8,2) -- (24,-4);
    \draw[-{Stealth[length=4.5pt]}, bkPurple, line width=1.4pt] (8,3) -- (30,3);
    \node[anchor=north, font=\sffamily\footnotesize, inner sep=0pt] at (16,-9) {矢印：2人で引く（4-10）};
  \end{scope}
\end{tikzpicture}\par}
\end{column}

\begin{taisetsu}
第4巻の道具は、図形（形）・三角関数（回転と波）・ベクトル（向き）の3つ。\par
どれも、「長さと角度で、場所や向きを数にする」ことから始まる。
\end{taisetsu}

\noindent 3つの部は、次のようにつながって、第5巻と物理シリーズ（力学）へ続いていきます。

\begin{bkfigure}
\begin{tikzpicture}[x=1mm, y=1mm,
  pb/.style={draw=bkTheme!80!black, line width=1pt, fill=white, rounded corners=2mm, minimum width=38mm, minimum height=27mm, inner sep=0pt},
  ph/.style={fill=bkTheme, text=bkOn, rounded corners=1mm, font=\sffamily\bfseries\footnotesize, inner xsep=1.5mm, inner ysep=0.8mm},
  pt/.style={font=\footnotesize, align=left, text width=34mm, inner sep=0pt},
  nb/.style={draw=bkNavy, line width=1pt, fill=white, rounded corners=2mm, minimum width=56mm, minimum height=10mm, font=\sffamily\small, align=center}]
  \node[pb] (p1) at (21,40) {};
  \node[pb] (p2) at (66,40) {};
  \node[pb] (p3) at (111,40) {};
  \node[ph, anchor=north west] at (3.5,51) {第1部 図形};
  \node[ph, anchor=north west] at (48.5,51) {第2部 三角関数};
  \node[ph, anchor=north west] at (93.5,51) {第3部 ベクトル};
  \node[pt, anchor=north west] at (4,45) {・証明（4-1）\\・ラジアン（4-2）\\・三平方の定理（4-4）};
  \node[pt, anchor=north west] at (49,45) {・sin・cos（4-6）\\・余弦定理（4-7）\\・回る点の影（4-8）};
  \node[pt, anchor=north west] at (94,45) {・矢印のたし算（4-10）\\・内積と仕事（4-12）\\・外積と回す力（4-14）};
  \draw[bk link] (p1) -- (p2);
  \draw[bk link] (p2) -- (p3);
  \node[nb] (v5) at (33,4) {第5巻\enspace 微分と積分};
  \node[nb] (ph) at (99,4) {物理シリーズ（力学）};
  \draw[bk link] (p2.south) -- (v5.north);
  \draw[bk link] (p1.south) -- (v5.north);
  \draw[bk link] (p3.south) -- (ph.north);
  \draw[bk link] (p2.south) -- (ph.north);
\end{tikzpicture}
\bkfigcaption{第4巻の道のり。矢印の先で、この巻の道具を使う。}
\end{bkfigure}

\clearpage
%% ---------------- 物理（力学）とのつながり ----------------
\bksection{butsuri}{このあと学ぶ物理（力学）とのつながり}

\label{bk:4-0:butsuri}この巻のすぐあとには、物理シリーズの「力学」が続きます。力学では、物の動きと力を、式で説明します。そこで使う道具の多くが、この巻で手に入ります（下の表の「使う道具」は、いまはわからなくてだいじょうぶ）。

\par\addvspace{4pt}\noindent\begin{minipage}[c]{0.44\linewidth}\centering
\begin{tikzpicture}[x=1mm, y=1mm]
  % 斜面（30°）と箱、重力と、その2つの分かれ
  \fill[bkGray!20] (0,0) -- (52,0) -- (52,30.02) -- cycle;
  \draw[bkGray!80!black, line width=0.9pt] (0,0) -- (52,0) -- (52,30.02) -- cycle;
  \begin{scope}[shift={(32,18.48)}, rotate=30]
    \filldraw[fill=bkOrange!45, draw=bkOrange!70!black, line width=0.6pt] (-5,0) rectangle (5,8);
  \end{scope}
  \coordinate (G) at ($(32,18.48)+(120:4)$);
  \draw[-{Stealth[length=4pt]}, bkInk, line width=1.2pt] (G) -- ++(0,-16);
  \draw[-{Stealth[length=4pt]}, bkBlue, line width=1.1pt] (G) -- ++(210:8);
  \draw[-{Stealth[length=4pt]}, bkRed, line width=1.1pt] (G) -- ++(300:13.86);
  \node[anchor=east, font=\sffamily\footnotesize, text=bkBlue!80!black, inner sep=1pt] at ($(G)+(210:8.5)$) {坂にそう向き};
  \node[anchor=north, font=\sffamily\footnotesize, inner sep=1pt] at ($(G)+(0,-16.5)$) {重力};
  \draw[bkInk, line width=0.5pt] (8,0) arc (0:30:8);
  \node[font=\sffamily\footnotesize, inner sep=0.5pt] at (15:11.5) {30°};
\end{tikzpicture}
\end{minipage}\hfill
\begin{minipage}[c]{0.53\linewidth}
たとえば、坂の上に置いた箱にはたらく重力（黒）は、「坂にそってすべらせようとする力」（青）と「坂をおす力」（赤）の2つに分けて考えます。分けた力の大きさは、坂の角度から $\sin$・$\cos$ で計算できます（4-6）。\enspace\bkivakore
\end{minipage}

\begin{bkcenter}
\begin{bktable}{colspec={Q[l,m,wd=58mm] Q[l,m,wd=38mm] Q[c,m,wd=21mm]}, rowsep=0.5pt, colsep=4pt, cells={font=\small}, row{1}={font=\sffamily\bfseries\small}}
力学でやること & 使う道具 & 章 \\
坂の上の物にはたらく力を分ける & 三角比（$\sin$・$\cos$） & 4-6 \\
2つの力を合わせた力の大きさ & 余弦定理 & 4-7 \\
3つの力のつり合い & 正弦定理・成分 & 4-7・4-11 \\
川を渡る船、電車の中から見た動き & ベクトルの和と差 & 4-10・4-11 \\
ボールの速度を、たてと横に分ける & ベクトルの成分 & 4-11 \\
物の重心 & 位置ベクトルの平均 & 4-11 \\
力がした仕事とエネルギー & 内積 & 4-12 \\
円運動の速さと周期 & ラジアン、角速度 & 4-2・4-8 \\
ボールがいちばん遠くへ飛ぶ角度 & 2倍角の式 & 4-9 \\
物を回す力（力のモーメント） & 外積 & 4-14 \\
惑星の動き & 極座標、だ円 & 4-8・4-15 \\
\end{bktable}
\bktabcaption{力学で使う、第4巻の道具}
\end{bkcenter}

\clearpage
%% ---------------- ⑤ 新しい記号 ----------------
\bksection{card}{この巻で出会う新しい記号}

第4巻では、次のような新しい記号に出会います。いまは意味がわからなくてもだいじょうぶです。どの記号も、使う前に場面と図でていねいに説明します（記号カードは巻末にも並べてあります）。

\begin{bkcenter}
\begin{bktable}{colspec={Q[c,m,wd=24mm] Q[l,m,wd=40mm] Q[l,m,wd=50mm] Q[c,m,wd=10mm]}, rowsep=0.3pt, colsep=3pt, cells={font=\small}, row{1}={font=\sffamily\bfseries\small}}
記号 & 読み方 & ひとことで言うと & 章 \\
$\angle\mathrm{ABC}$ & 角 ABC & 点Bにできる角 & 4-1 \\
$\triangle\mathrm{ABC}$ & 三角形 ABC & 3つの点A・B・Cを結んだ三角形 & 4-1 \\
$\perp$、$\parallel$ & 垂直、平行 & 直角に交わる、交わらない & 4-1 \\
$\theta$ & シータ & 角の大きさ & 4-2 \\
rad & ラジアン & 角の単位（弧の長さ÷半径） & 4-2 \\
$(r,\ \theta)$ & アール、シータ & 距離と向きの角で表した場所（極座標） & 4-2 \\
$\sojisym$ & 相似 & 形が同じ（大きさはちがってよい） & 4-3 \\
$\equiv$ & 合同 & ぴったり重なる & 4-3 \\
$\sin$・$\cos$・$\tan$ & サイン・コサイン・タンジェント & 直角三角形の辺の比 & 4-6 \\
$\omega$ & オメガ & 角速度。1秒に回る角 & 4-8 \\
$\vec{a}$、$|\vec{a}|$ & ベクトル $a$、その大きさ & 向きと大きさをもつ量 & 4-10 \\
$\vec{a}\cdot\vec{b}$ & $a$ ドット $b$ & 内積。どれだけ同じ向きか & 4-12 \\
$\boldsymbol{a}$、$\vec{i}$・$\vec{j}$・$\vec{k}$ & 太字の $a$、$i$・$j$・$k$ & 専門書のベクトルの書き方、軸の向きの矢印 & 4-13 \\
$\vec{a}\times\vec{b}$ & $a$ クロス $b$ & 外積。回す力の計算 & 4-14 \\
\end{bktable}
\bktabcaption{この巻で出会う記号（くわしくはそれぞれの章で）}
\end{bkcenter}

同じ形の記号が、ちがう意味で使われることがあります。$\triangle$（三角形）は、2-3の $\Delta$（デルタ、変化量）とは別の記号です。$\vec{a}\times\vec{b}$ の「×」は、ふつうの数のかけ算とはちがう計算です（4-14）。第5巻で使う記号（$\lim$ など）は、この巻では使いません。

%% ---------------- この巻の読み方 ----------------
\bksection{hon}{この巻の読み方}

\label{bk:4-0:yomikata}
\begin{bkfigure}
\begin{tikzpicture}[x=1mm, y=1mm]
  \bkivatalk{0}{sora}{ソラ}{第3巻の最後に、デルタが「第4巻からは、漢字にふりがながつかないよ」って言ってたね。}
  \bkivatalk{-20}{mina}{ミナ}{読めない漢字が出てきたら、どうしよう……。図形や三角関数も、むずかしそう。}
  \bkivatalk{-44}{delta}{デルタ}{だいじょうぶ。新しいことばには、読みを書くよ。図形も三角関数も、まず図をかいて、はかって、気づいてから式にするよ。わからなくなったら、章番号の札の章にもどろう。}
\end{tikzpicture}
\end{bkfigure}

\begin{bkfigure}
\begin{tikzpicture}[x=1mm, y=1mm]
  \bkivacard{(0,0)}{omoidasu}{bkGray}{①\enspace ふりがなのかわりに}
    {新しいことばには、はじめて出てきた所で、読みを（　）で書きます。\\[2pt]
     \quad 例：弧（こ）、錯角（さっかく）\\[2pt]
     声に出して読んでみよう。}
  \bkivacard{(68,0)}{reidai}{bkBlue}{②\enspace「式」と「答え」}
    {例題の中の札が、ひらがなの\\
     \quad\bkivapill{しき}\enspace\bkivapill{こたえ}\enspace から、漢字の\\
     \quad\bkivapill{式}\enspace\bkivapill{答え}\\
     に変わります。中身は同じです。}
  \bkivacard{(0,-52)}{doguu}{bkOrange}{③\enspace 図をかく道具}
    {\bkivaitem{定規・分度器・コンパスを用意する。}\bkivaitem{まず図をノートに大きくかき、わかった長さや角を書きこむ。}}
  \bkivacard{(68,-52)}{te}{bkNavy}{④\enspace 電卓と表計算}
    {\bkivaitem{電卓・表計算・プログラムを使ってよい。}\bkivaitem{関数電卓の DEG（度）と RAD（ラジアン）の切りかえ（4-2）を、おす前に確かめる。}}
\end{tikzpicture}
\end{bkfigure}
